\documentclass[11pt]{article}
\usepackage[a4paper,margin=1in]{geometry}
\usepackage{amsmath,amssymb}
\usepackage{booktabs}
\usepackage{hyperref}
\title{Cost-Weighted GEC for DELCODE Whole-Brain\\Method and Implementation Summary}
\author{Project Notes}
\date{\today}

\begin{document}
\maketitle

\section*{Goal}
This document summarizes the cost-weighted objective used for converter classification in the GEC pipeline and how it is applied in the whole-brain cross-validation training notebook.

\section*{Task Definition}
The model performs binary graph-level classification:
\begin{itemize}
\item Class $0$: non-converter
\item Class $1$: converter
\end{itemize}
For each graph sample $i$, the model outputs a logit $z_i \in \mathbb{R}$ and label $y_i \in \{0,1\}$.

\section*{Base BCEWithLogits Loss}
The unreduced per-sample binary cross-entropy with positive weighting is:
\begin{equation}
\ell_i = -\Big(\alpha\, y_i\log\sigma(z_i) + (1-y_i)\log\big(1-\sigma(z_i)\big)\Big),
\end{equation}
where $\sigma(\cdot)$ is the sigmoid function and
\begin{equation}
\alpha = \texttt{pos\_weight} = \frac{N_{-}}{N_{+}}.
\end{equation}
$N_{+}$ and $N_{-}$ are the positive (converter) and negative (non-converter) sample counts in the current training fold.

\section*{Class Cost Weights}
A second class-balancing term is computed from fold label counts:
\begin{equation}
 w_0 = \frac{N}{2N_{-}}, \qquad w_1 = \frac{N}{2N_{+}},
\end{equation}
with $N = N_{+}+N_{-}$. In the implementation, these are optionally normalized by their mean:
\begin{equation}
\tilde{w}_c = \frac{w_c}{(w_0+w_1)/2}, \quad c\in\{0,1\}.
\end{equation}
Per-sample cost weight is then assigned by class:
\begin{equation}
\omega_i =
\begin{cases}
\tilde{w}_1, & y_i=1 \\
\tilde{w}_0, & y_i=0.
\end{cases}
\end{equation}

\section*{Weighted Cost Averaging}
The final training loss is the weighted average of unreduced BCE losses:
\begin{equation}
\mathcal{L}_{\text{CW}} = \frac{\sum_{i=1}^{B} \omega_i\,\ell_i}{\sum_{i=1}^{B} \omega_i + \varepsilon},
\end{equation}
where $B$ is batch size and $\varepsilon$ is a small constant for numerical stability.

\section*{How It Was Applied in This Project}
In the cost-weighted whole-brain notebook, each CV fold does the following:
\begin{enumerate}
\item Build train/validation subsets from subject-indexed scan graphs.
\item Compute \texttt{pos\_weight} from fold labels.
\item Compute \texttt{class\_cost\_weights} from fold labels.
\item Train with
\begin{itemize}
\item unreduced BCEWithLogits (with \texttt{pos\_weight}), then
\item weighted cost averaging using \texttt{class\_cost\_weights}.
\end{itemize}
\item Evaluate fold AUC, sensitivity, specificity, and F1 with threshold selected by Youden's $J$ statistic:
\begin{equation}
J = \text{TPR} - \text{FPR}.
\end{equation}
\end{enumerate}

\section*{Interpretation}
\begin{itemize}
\item If converters are minority in a fold, $\tilde{w}_1 > \tilde{w}_0$ and converter errors contribute more.
\item If non-converters are minority, $\tilde{w}_0 > \tilde{w}_1$.
\item Using both \texttt{pos\_weight} and class-cost averaging applies balancing at two levels:
\begin{itemize}
\item inside BCE positive term scaling,
\item outside BCE as weighted batch averaging.
\end{itemize}
\end{itemize}

\section*{Implementation Locations}
\begin{itemize}
\item Loss definition and weighted averaging: \texttt{MODEL/model/CostWeightedGEC/train.py}
\item Weight computation: \texttt{MODEL/model/CostWeightedGEC/utils.py}
\item Fold-wise use in training: \texttt{MODEL/notebooks/COST\_WEIGHTED\_GEC\_DELCODE\_WHOLE\_BRAIN.ipynb}
\end{itemize}

\end{document}
